\documentclass[10pt]{article}

% ---------- Document Identity (update these only) ----------
\newcommand{\DocStage}{}
\newcommand{\DocTitle}{Your Road to Tier One SOF}
\newcommand{\ProgramName}{182-Week Civilian-to-Selection Readiness Pipeline}
\newcommand{\ProgramSubtitle}{}

% ---------- Page ----------
\usepackage[legalpaper,left=0.5in,right=0.5in,top=0.6in,bottom=0.45in]{geometry}

% ---------- Typography ----------
\usepackage[T1]{fontenc}
\usepackage[utf8]{inputenc}
\usepackage{libertinus}
\usepackage{microtype}
\usepackage{parskip}
\usepackage{ragged2e}
\usepackage{textcomp}
\AtBeginDocument{\color{black}}
\AtBeginDocument{\pagecolor{white}}
\AtBeginDocument{\justifying}

% ---------- Landscape pages ----------
\usepackage{pdflscape}

% ---------- Color ----------
\usepackage[table]{xcolor}
\definecolor{StageBlue}{HTML}{1F4E79}
\definecolor{StageBlueDark}{HTML}{163A5A}
\definecolor{LightGray}{HTML}{F2F4F7}
\definecolor{MidGray}{HTML}{D9DEE5}
\definecolor{RuleGray}{HTML}{C7CED8}
\definecolor{HeaderInk}{HTML}{000000}
\definecolor{Ink}{HTML}{000000}

% ---------- Utility ----------
\usepackage{enumitem}
\sloppy
\setlength{\emergencystretch}{2em}

% ---------- Boxes / UI ----------
\usepackage[most]{tcolorbox}

\tcbset{
  charterTitle/.style={
    enhanced,
    colback=StageBlue!4,
    colframe=StageBlue,
    boxrule=1.25pt,
    arc=3pt,
    left=16pt,right=16pt,top=14pt,bottom=14pt,
    borderline north={3.2pt}{0pt}{StageBlue},
    borderline west={4pt}{0pt}{StageBlue},
    borderline south={1.25pt}{0pt}{StageBlue},
    drop shadow={black!25!white},
  },
  charterRule/.style={
    enhanced,
    colback=LightGray,
    colframe=RuleGray,
    boxrule=0.85pt,
    arc=3pt,
    left=12pt,right=12pt,top=10pt,bottom=10pt,
    borderline west={3pt}{0pt}{StageBlue},
  },
  charterBadge/.style={
    enhanced,
    colback=StageBlue,
    coltext=white,
    colframe=StageBlueDark,
    boxrule=0.6pt,
    arc=2pt,
    left=6pt,right=6pt,top=3pt,bottom=3pt,
  }
}
\newtcbox{\Badge}{nobeforeafter,charterBadge}

% ---------- Lists ----------
\setlist[itemize]{leftmargin=1.5em, itemsep=0.25em, topsep=0.25em}
\setlist[enumerate]{leftmargin=1.8em, itemsep=0.25em, topsep=0.25em}

% ---------- TOC ----------
\usepackage{tocloft}
\setcounter{tocdepth}{3}
\setlength{\cftbeforesecskip}{3pt}
\setlength{\cftbeforesubsecskip}{2pt}
\setlength{\cftbeforesubsubsecskip}{0pt}
\renewcommand{\cftsecfont}{\bfseries}
\renewcommand{\cftsecpagefont}{\bfseries}
\renewcommand{\cftsubsecfont}{\normalfont}
\renewcommand{\cftsubsecpagefont}{\normalfont}
\renewcommand{\cftsubsubsecfont}{\normalfont}
\renewcommand{\cftsubsubsecpagefont}{\normalfont}
\setlength{\cftsecindent}{0pt}
\setlength{\cftsubsecindent}{1.25em}
\setlength{\cftsubsubsecindent}{2.8em}
\setlength{\cftsecnumwidth}{2.1em}
\setlength{\cftsubsecnumwidth}{2.8em}
\setlength{\cftsubsubsecnumwidth}{3.6em}

% Remove the built-in TOC heading so only the custom heading is shown
\makeatletter
\renewcommand{\tableofcontents}{\@starttoc{toc}}
\makeatother

% ---------- Header/Footer: page number only ----------
\usepackage{fancyhdr}
\pagestyle{fancy}
\fancyhf{}
\renewcommand{\headrulewidth}{0pt}
\renewcommand{\footrulewidth}{0pt}
\fancyfoot[C]{\thepage}

% ---------- Section formatting ----------
\usepackage{titlesec}

\titleformat{\section}
  {\Large\bfseries\color{Ink}}
  {\thesection}{0.6em}{}
  [\vspace{0.25em}\textcolor{StageBlue}{\rule{\linewidth}{0.8pt}}]

\titleformat{\subsection}
  {\large\bfseries\color{Ink}}
  {\thesubsection}{0.6em}{}

\titleformat{\subsubsection}
  {\normalsize\bfseries\color{Ink}}
  {\thesubsubsection}{0.6em}{}

\titlespacing*{\section}{0pt}{0.95em}{0.35em}
\titlespacing*{\subsection}{0pt}{0.65em}{0.25em}
\titlespacing*{\subsubsection}{0pt}{0.55em}{0.2em}

% ---------- Table engine ----------
\usepackage{tabularray}
\UseTblrLibrary{booktabs}

% ---------- tabularray: GLOBAL table treatment ----------
\SetTblrInner{
  cells = {fg=black, bg=white, valign=t},
}

% ---------- Header helpers ----------
\newcommand{\Hdr}[1]{\shortstack[c]{\strut #1\strut}}

% ---------- Lines ----------
\newcommand{\TitleLine}{\textcolor{StageBlue}{\rule{0.98\linewidth}{0.95pt}}}
\newcommand{\SubTitleLine}{\textcolor{RuleGray}{\rule{0.98\linewidth}{0.65pt}}}

% ---------- Continued on next page ----------
\newcommand{\ContinuedOnNextPageInline}{%
  \par
  \vfill
  \noindent\textcolor{RuleGray}{\rule{\linewidth}{1pt}}\vspace{-2pt}
  \noindent\hfill\textit{Continued on next page}\par\vspace{-10pt}
  \noindent\textcolor{RuleGray}{\rule{\linewidth}{1pt}}
  \clearpage
}

% ---------- PDF hyperlinks ----------
\usepackage[hidelinks]{hyperref}
\hypersetup{
  colorlinks=false,
  pdfborder={0 0 0},
  linktoc=all,
  pdftitle={\DocTitle},
  pdfauthor={\ProgramName}
}

% =========================================================
\begin{document}

\begin{tcolorbox}
\begin{center}
{\LARGE\bfseries Stage 4 — Phase 00 Specification}\\[0.35em]
{\Large\bfseries SOF Physical Development Transformation Program}\\[0.25em]
{\large 182-Week Civilian-to-Selection Readiness Pipeline}\\[0.65em]
\TitleLine
\end{center}
\end{tcolorbox}

\vspace{0.35em}

\section*{Stage 4 — Phase 00 Specification}
\phantomsection
\addcontentsline{toc}{section}{Stage 4 — Phase 00 Specification}

\subsection*{Introduction}
\phantomsection
\addcontentsline{toc}{subsection}{Introduction}

\subsubsection*{1) Purpose}

\begin{itemize}
\item Stage 4 exists to specify Phase 00 execution doctrine that is safe, repeatable, and audit-ready under an obese, deconditioned start-state with elevated musculoskeletal and cardiometabolic risk.
\item Stage 4 establishes the minimum deployable operating system for Phase 00 under start-from-zero constraints, ensuring Phase 00 can be executed with stable documentation, stable controls, and enforceable safety posture.
\item Stage 4 is the implementation bridge between upstream governance and downstream phase design, equipment timing, and category manuals, without redefining or weakening any upstream doctrine.
\item Stage 4 defines Phase 00 only and constrains how Phase 00 produces admissible evidence, how it handles deviations, and how it establishes safe entry posture for Phase 01.
\item Stage 4 contains and governs these Phase 00 deliverables as a closed set:
  \begin{itemize}
  \item 4.A — Phase 00 Overview — Phase Row Template
  \item 4.B — Phase 00 Timeline and Microcycle Structure
  \item 4.C — Movement Menu, Regressions, and Progressions
  \item 4.D — Intensity, Volume, and Progression Rules
  \item 4.E — Testing Within Phase 00 — Scaled Use of Stage 2 Protocols
  \item 4.F — Safety, Red Flag Criteria, and Stop Rules for Phase 00
  \item 4.G — Exit Standards for Phase 00 — Safe Entry to Phase 01
  \end{itemize}
\end{itemize}

\subsubsection*{2) Scope}

\paragraph*{2.1 Inclusions}

\begin{itemize}
\item Phase 00 execution boundaries and intent, including conservative ramp and baseline creation under low-impact bias and high-risk novice posture.
\item Phase 00 session structure components at the element level, aligned to Stage 1.F validity discipline, admissibility posture, and documentation minimums.
\item Movement menu and scaling boundaries suitable for an obese beginner posture, including regressions, progressions, and non-negotiable regression triggers.
\item Phase 00 testing posture executed only inside the Phase 00 protected Deload+Test windows defined upstream and locked in Stage 3 and Stage 4.A.
\item Scaled use posture for measurement: where a Phase 00 test construct is already defined in Stage 2, Stage 2 protocol and logging conventions govern admissibility; where a Phase 00 construct is not yet defined in Stage 2, it must be treated as a Phase 00-local construct and flagged for reconciliation rather than treated as Stage 2-backed.
\item Safety overlays, stop triggers, and escalation posture as applied specifically in Phase 00, implemented as scheduling and execution controls without medical diagnosis or treatment guidance.
\item Phase 00 exit standards and gate decision posture aligned to Stage 1.F gate outcomes and Stage 3 reslot discipline.
\end{itemize}

\paragraph*{2.2 Exclusions}

\begin{itemize}
\item Any new tests, thresholds, domains, or identifiers.
\item Any rewriting of Stage 2 protocol cards, Stage 2 metrics inventory, or Stage 2 logging conventions.
\item Any rewriting of Stage 3 test calendar placement, Deload+Test doctrine, stacking doctrine, spacing doctrine, or reslot doctrine.
\item Any Phase 01–Phase 13 programming, phase purposes, or macro design content.
\item Any equipment category expansion, retailer guidance, or acquisition guidance.
\item Any clinical diagnosis, treatment advice, or prescription guidance.
\end{itemize}

\subsubsection*{3) How to Use}

\paragraph*{Coach Admin Use}

\begin{itemize}
\item Use Stage 4 as the Phase 00 operating manual that constrains execution and prevents drift into unsafe exposure escalation or undocumented substitutions.
\item Enforce Stage 0 constraints and Stage 1.F validity doctrine as non-negotiable controls in every session and every test window.
\item Align Phase 00 testing to Stage 3 protected Deload+Test windows and to Stage 2 protocolized constructs where they exist, without rewriting Stage 2 procedures or standards.
\item Enforce admissibility posture: gate decisions rely on admissible evidence only; incomplete logs, missing lock fields, and compromised conditions are treated as inadmissible for gate use.
\item Apply substitution and reslot posture when environment, access, or safety conditions change; preserve safety and validity rather than preserving schedule.
\item Treat any Phase 00 test construct that lacks Stage 2 protocol and metric identifiers as a Phase 00-local construct and flag it as Requires Stage 8 review; do not treat such constructs as crosswalk join keys or as Stage 2-backed gate evidence until reconciled upstream.
\item Use Phase 00 documentation to produce gate-grade evidence for Phase 00 exit decisions by enforcing same-day logging, stable identifiers where defined, and explicit deviation recording.
\end{itemize}

\paragraph*{Trainee Use Under Coach Supervision}

\begin{itemize}
\item Phase 00 is reactivation and baseline creation under conservative risk posture; it is not a performance phase and is not a phase for “testing your limits.”
\item Log daily and per session to the minimum dataset required for admissibility posture, including required session elements, recovery markers, and foot and pain flags.
\item Follow non-negotiable safety posture: red-flag symptoms and mechanics-altering pain trigger \textbf{STOP} \textbf{NOW} and escalation posture as directed by governance.
\item Maintain no novelty near Deload+Test windows: do not change footwear systems, route classes, tool methods, or measurement standards inside decision windows unless explicitly documented as a comparability break.
\item Accept reslot discipline: if conditions are unsafe or validity is compromised, the correct action is to reslot, not to force an attempt.
\end{itemize}

\subsubsection*{4) Inputs and Assumptions}

\paragraph*{Inputs}

\begin{itemize}
\item Start state and constraints from Stage 0, including start-from-zero posture, operating constraints, safety boundaries, and end-state envelope framing.
\item Governance and validity doctrine from Stage 1.F, including validity tags, admissibility posture, stop rules, gate outcomes, and change control posture.
\item Stage 2 domain taxonomy, metric inventory, protocol cards, tier thresholds, and Stage 2 logging conventions, as the authoritative measurement layer wherever Phase 00 uses Stage 2-defined constructs.
\item Stage 3 protected Deload+Test windows and adjustment doctrine, as the authoritative scheduling layer for when Phase 00 may attempt gate-grade test evidence.
\item Phase 00 duration and window placement as defined upstream, including the locked Phase 00 Deload+Test windows in Stage 3 and the Phase 00 batteries and exit posture in Stage 4.A.
\item Variable access posture consistent with start-from-zero constraints, including environment volatility and indoor substitution posture when outdoor conditions are unsafe.
\end{itemize}

\paragraph*{Assumption Ledger: not required for Stage 4 Intro.}

\subsubsection*{5) Interfaces}

\paragraph*{Upstream Dependencies}

\begin{itemize}
\item Stage 0 identity, constraints, end-state envelope posture, and safety posture, including environmental safety dominance.
\item Stage 1.F governance: session validity, test validity, admissibility, stop posture, gate outcomes, freeze windows, and change control.
\item Stage 2 measurement architecture: domain taxonomy, metric inventory, protocol cards, thresholds, and logging conventions.
\item Stage 3 scheduling architecture: Deload+Test doctrine, protected windows, stacking and spacing doctrine, reslot posture, and adjustment rules.
\end{itemize}

\paragraph*{Downstream Consumers}

\begin{itemize}
\item Stage 5 Phase 01–13 macro design alignment, which consumes Phase 00 exit posture as the safe entry boundary for Phase 01.
\item Stage 6 equipment and resource timing, ensuring Phase 00 does not assume capabilities not yet introduced and that Phase 00 prerequisites are feasible.
\item Stage 7 category overviews, which operationalize tools and systems used during Phase 00 without expanding categories.
\item Stage 9 packaging and binder-based audit traceability, which requires Phase 00 records to remain decision-grade and crosswalk-ready.
\end{itemize}

\paragraph*{Crosswalk Hooks}

\begin{itemize}
\item Phase identifiers Phase 00–Phase 13.
\item Phase Week and Session numbering conventions.
\item Validity tags and gate outcomes defined upstream.
\item Protocol Card IDs \textbf{C1}-\textbf{C18} and Metric IDs \textbf{MET-2B}-### as stable join keys for any protocolized Phase 00 constructs.
\item Environment unit identifiers and tooling consistency posture where Stage 2.E defines them.
\item Requires Stage 8 review: any Phase 00 construct referenced in Stage 4 that is not already defined in Stage 2 with a stable Protocol Card \textbf{ID} and Metric \textbf{ID} must be reconciled upstream before it can be treated as crosswalk join-key evidence.
\end{itemize}

\subsubsection*{6) Gate and Acceptance Criteria}

\begin{itemize}
\item Pass if Stage 4 Intro accurately frames Phase 00 and the deliverable set 4.A–4.G as a closed set without redefining upstream doctrine.
\item Pass if safety governance and stop posture are explicitly anchored and do not conflict with Stage 0 and Stage 1.F authorities.
\item Pass if validity and logging posture references Stage 1.F and Stage 2 logging conventions without alteration and preserves admissibility posture for gate decisions.
\item Pass if testing posture is explicitly constrained to Stage 3 protected windows and does not imply work-week gate-grade testing.
\item Pass if any mismatch between Phase 00 test constructs and Stage 2 identifiers is explicitly flagged as Requires Stage 8 review and not silently corrected or implied as already resolved.
\item Fail if Phase 00 is framed as a performance phase, or if the Intro implies unsafe early exposure escalation, or if it implies underwater or breath-hold exposure.
\item Fail if Stage 4 claims Stage 2 coverage for Phase 00 constructs that are not actually defined in Stage 2, without flagging Requires Stage 8 review.
\end{itemize}

\subsubsection*{7) Common Failure Modes and Controls}

\begin{itemize}
\item Aggressive ramp beyond start-state tolerance → conservative progression boundaries, low-impact bias, and mandatory validity tagging and reslot posture.
\item Silent substitutions → mandatory \textbf{MODIFIED} tagging, reason-code posture, and explicit deviation notes; undocumented drift is inadmissible for gate use.
\item Novelty changes near Deload+Test windows → no-novelty posture, lock-field discipline, and reslot rather than forced testing.
\item Treating Phase 00 as a performance phase → explicit Phase 00 mission posture and exit standards framed as safe entry to Phase 01, not early peak outputs.
\item Missing logs or reconstructed logs → inadmissibility posture; missing-data rules; do not claim gate credit on incomplete evidence.
\item Unsafe environment exposure → substitution to safe indoor equivalents only when comparability is preserved and logged; otherwise reslot per Stage 3.
\item Ignoring pain that alters mechanics → \textbf{STOP} \textbf{NOW} posture for the provoking exposure; conservative downshift; document and reslot testing as required.
\item Phase 00 gate items not traceable to Stage 2 identifiers → flag Requires Stage 8 review; treat as Draft - Non-Deployable crosswalk evidence until reconciled; do not claim Stage 2-backed admissibility.
\item Equipment feasibility drift under start-from-zero posture → treat missing capabilities as constraints; apply \textbf{HOLD} posture and reslot test windows rather than improvising unlogged substitutions.
\item Evidence used outside protected windows → enforce that gate-grade evidence is produced only in Deload+Test windows; work-week attempts are trend-only and not gate-grade by default.
\end{itemize}

\subsubsection*{8) Versioning Notes}

\begin{itemize}
\item Stage 4 is part of Program v1.0 and is frozen except for cosmetic corrections that do not change meaning, enforcement posture, or interfaces.
\item Any substantive change requires controlled patch discipline consistent with governance, including change-log entry, validation check, and distribution control.
\item No mid-phase doctrinal edits without change control and crosswalk impact review.
\item Requires Stage 8 review remains the authoritative flag for unresolved mismatches between Phase 00 constructs and Stage 2 identifiers, and for any crosswalk integrity issues that must be corrected upstream before downstream artifacts can remain deployable.
\end{itemize}

\end{document}